\chapter{Lump Under Armpit (Axillary Mass)}

\section*{Definition}
A palpable swelling or mass in the axillary region, most commonly representing enlarged lymph nodes but also arising from breast tissue, skin lesions, or subcutaneous structures.

\section*{Pathophysiology}
Axillary lymph nodes drain the upper limb, breast, chest wall, and upper back. Lymphadenopathy results from infection (local or systemic), inflammatory conditions, or malignancy (lymphoma, breast cancer metastasis). Other causes include sebaceous cysts, lipomas, hidradenitis suppurativa, and accessory breast tissue.

\section*{Causes}
\begin{itemize}
  \item Lymphadenopathy (reactive from infection of arm/hand)
  \item Breast cancer metastasis
  \item Lymphoma or leukemia
  \item Sebaceous cyst or epidermoid cyst
  \item Lipoma (benign fatty tumor)
  \item Hidradenitis suppurativa (inflamed sweat gland)
  \item Cat scratch disease (Bartonella henselae)
  \item Tuberculosis or mycobacterial infection
  \item Sarcoidosis
  \item Vaccination reaction (recent COVID-19, flu, or other vaccine)
\end{itemize}

\section*{When to See a Doctor}
See your doctor if:
\begin{itemize}
  \item Severe or worsening symptoms
  \item Firm, fixed, non-tender lump, progressive enlargement, or lump with fever, night sweats, or unexplained weight loss
  \item Accompanied by other concerning signs
\end{itemize}

\section*{Related Symptoms}
\begin{itemize}
  \item Breast lump
  \item Arm swelling
  \item Fever
  \item Night sweats
  \item Weight loss
\end{itemize}
\textbf{Important:} If this symptom persists, your doctor will check for serious underlying causes.

\section*{Self-Care / Home Management}
\begin{itemize}
  \item Rest and avoid activities that make it worse.
  \item Maintain adequate hydration and a balanced diet.
  \item Over-the-counter remedies may provide symptomatic relief (consult a pharmacist or doctor).
  \item Monitor symptoms and seek professional advice if they persist or worsen.
  \item Avoid self-medicating with prescription medications without medical guidance.
\end{itemize}

\section*{How Your Doctor Will Evaluate You}
\begin{itemize}
  \item A thorough history and physical examination are the first steps in evaluation.
  \item Laboratory tests (blood work, urinalysis) may be ordered based on clinical suspicion.
  \item Imaging studies (X-ray, ultrasound, CT, MRI) may be indicated when structural pathology is suspected.
  \item Specialist referral is arranged when the diagnosis remains uncertain or requires specific expertise.
\end{itemize}

\section*{Treatment Options}
\begin{itemize}
  \item Treatment targets the underlying cause identified through diagnostic evaluation.
  \item Supportive care and symptom management are provided while awaiting definitive therapy.
  \item Medications, lifestyle modifications, and physical therapies may be prescribed as appropriate.
  \item Follow-up ensures treatment effectiveness and allows timely adjustments to the care plan.
\end{itemize}

\clearpage